\documentclass{article}
\begin{document}
\section*{Fubini under additive presentation changes at arbitrary finite shape}
The raw tag types, coefficient labels, binary tree shape, and
locations of multiplicative units stay fixed. At each leaf the
parent word and each continuation word can be reordered. The
comparison to the induced presentation is computed using the
actual augmented word transitions. At a leaf the read construction
is proved equal as an augmented arrow to the previously checked
literal raw native isomorphism.

On the separate route the vertical is the recursive independent
product of those isomorphisms. On the joint route it is the
accumulation of the recursively paired parent and continuation
reads, including their rectangular chart and augmentation laws.
These arrows are constructed independently of the Fubini comparison;
their square commutes for every finite binary tree, with internal
units permitted. The unit case uses terminality of the
multiplicative unit augmentation, not initiality of additive zero.
The theorem does not filter out empty fibres or empty families.

For one dependent input, reordering by composed parent and child
circuits gives the iterated ordered presentation with inverse
charts composed in reverse order. The native augmented transitions
obey the corresponding composition and identity laws after
transport at that explicit object equality. The proof compares
raw tags through the actual singleton cells, not their coefficient
labels. No identification of different circuits on a dormant word
with repeated raw tags is used.

This is an arbitrary-shape result in the native augmented
category. It does not package the indexed lift of arbitrary
independent structural paths, arbitrary mixed dependent
associativity pasting, or the separate independent audit of the
paper's ambient rig interface.
\end{document}
